\documentclass[10pt]{article}
\usepackage[a4paper,margin=22mm]{geometry}
\usepackage[T1]{fontenc}
\usepackage{lmodern}
\usepackage{amsmath,amssymb,bm,booktabs,array}
\usepackage[hidelinks]{hyperref}
\newcommand{\R}{\mathbb R}
\newcommand{\PP}{\mathbb P}
\newcommand{\dd}{\mathrm d}
\title{Active Flow Control from an Implicit Differential Representation}
\author{Mohamed Abdelghani}
\date{Numerical prototype and experiment protocol --- 7 November 2026}
\begin{document}
\maketitle

\paragraph{Scope.}
This section accompanies the active-flow-control implementation in the
simulation package. It describes a newly implemented numerical prototype;
its records are independent of the numerical tables in the earlier paper.
Published airfoil-control studies provide the physical motivation:
Garcia et al.\ use surface blowing and suction for a two-dimensional
NACA~0012 airfoil at $Re=3000$ \cite{Garcia2025Airfoil}, while
Montal\`a et al.\ study a three-dimensional separated wing at $Re=1000$
\cite{Montala2025Wing}. The actuator, sensing, and validation scope of the
present computation are stated explicitly below.

\section{Flow system and implicit differential}
Let $\bm v(\bm x,t)$ be an incompressible velocity field, $p$ pressure,
and $r$ the actuator command. A fixed, smooth body-force profile
$\mathcal B(\bm x)$ produces $\beta\mathcal B r$. With a stationary immersed
airfoil and an outer-domain relaxation layer, the continuum model behind
the projected computation is
\begin{align}
\mathcal R_m &={\bm v}_t+
 c_a(\bm v\cdot\nabla)\bm v+\nabla p-
 \nu\Delta\bm v+\kappa\chi\bm v
 +\sigma(\bm x)(\bm v-\bm v_\infty)-\beta\mathcal B r=0,
 \label{eq:afc_ns}\\
\mathcal R_d&=\nabla\cdot\bm v=0.
 \label{eq:afc_div}
\end{align}
Here $\chi$ is the smoothed obstacle mask, $\kappa>0$ is the
penalization strength, and $\sigma\geq0$ is the prescribed relaxation
coefficient; the nominal plant has $c_a=1$ and $\nu=Re^{-1}$.
Penalization is motivated by the approach of
Angot et al.\ \cite{Angot1999Penalization}; velocity projection follows
the incompressible-flow construction of Chorin \cite{Chorin1968Projection}.
The airfoil geometry belongs to the NACA four-digit family
\cite{Jacobs1933NACA}. These references motivate numerical ingredients;
they do not validate the spatial or temporal resolution used here.

With $a=\bm v_t$ and $z=\operatorname{col}(\bm v,a,p,r)$, the actual
Fr\'echet differential of the momentum residual is
\begin{equation}
 D\mathcal R_m(z)\,\delta z
 =\mathcal K(\bm v)\delta\bm v+\delta a+
 \nabla\delta p-\beta\mathcal B\delta r,
 \label{eq:afc_frechet}
\end{equation}
where
\begin{equation}
\mathcal K(\bm v)\delta\bm v
 =c_a\bigl((\delta\bm v\cdot\nabla)\bm v+
  (\bm v\cdot\nabla)\delta\bm v\bigr)
  -\nu\Delta\delta\bm v+
  (\kappa\chi+\sigma)\delta\bm v.
\end{equation}
The divergence row contributes $\nabla\cdot\delta\bm v$.
Thus $F=D_z\operatorname{col}(\mathcal R_m,\mathcal R_d)$ is a
state-dependent differential operator. Its evaluated entries change as
the flow evolves even when the physical coefficients are constant.
Within a regime of constant physical coefficients, a sufficiently smooth
admissible trajectory satisfies
\begin{equation}
 F(z)\dot z=0.
 \label{eq:afc_compatibility}
\end{equation}
This relation involves $\bm v_{tt}$ and $p_t$. Differentiation of sampled
data must therefore be specified rather than treating these rates as exact
experimental measurements.

\section{Learning and a state-dependent output readout}
The executable plant eliminates pressure using the discrete solenoidal
Fourier projection $\PP_h$, including its declared spectral cutoff. Let
$X$ contain both velocity components, $A=\dot X$, and define
\begin{align}
 N_h(X)&=\PP_h[(X\cdot\nabla_h)X],\qquad
 L_h(X)=\Delta_hX,\qquad B_h=\PP_h\mathcal B,\\
 K_h(X)&=\PP_h[\sigma\bm v_\infty-(\kappa\chi+\sigma)X].
\end{align}
The simulator evolves $\dot X=-c_aN_h(X)+\nu L_h(X)+\beta B_hr+K_h(X)$.
The learned implicit family and its actual partials are
\begin{align}
 f_\theta(X,A,r)&=A+c_aN_h(X)-\nu L_h(X)-\beta B_hr-K_h(X),
 \quad \theta=(c_a,\nu,\beta),\label{eq:afc_discrete_residual}\\
 F_\theta&=\left[
 c_aDN_h(X)-\nu L_h-DK_h\quad I\quad -\beta B_h
 \right].\label{eq:afc_discrete_partials}
\end{align}
Within each constant-coefficient regime the physical relation is
$f_{\theta_\star}(X,\dot X,r)=0$.
For example,
$DN_h(X)\delta X=\PP_h[(\delta X\cdot\nabla_h)X+(X\cdot\nabla_h)\delta X]$.
The code evaluates the convection term in its equivalent projected rotational
form. Only three coefficients of this specified operator family are learned;
the geometry, differential operators, and forcing profiles are supplied.
The identity $A$ block excludes the trivial zero representation. The
prescribed residual family fixes its zero level; no free integration
constant is identified by differentiating observations. Learning uses
observed variations through the incompatibility
$\xi_\theta=F_\theta(z)\dot z$, with the candidate parameters held fixed
when taking the context derivative. The continuous counterpart of the
zero-initialized exponential memory is
\begin{equation}
 \mathcal E_t(\theta)=\frac{1}{2\tau_m}
  \int_{t_0}^{t}e^{-(t-s)/\tau_m}
  \|F_\theta(z(s))\dot z(s)\|_h^2\,\dd s,
 \label{eq:afc_energy}
\end{equation}
where $\|v\|_h^2$ is the mean squared value over both velocity components
and every grid point. The memory is not divided by an accumulated
weight normalization. The feasible
parameter set is the declared box
$[0.2,2]\times[0.0005,0.05]\times[0.15,2]$, which also keeps the actuator
gain positive. The sampled update is specified below.

Small training incompatibility alone is insufficient evidence of successful
identification. We also evaluate a partial-action error, the actuator-partial
error, and the next-interval conditional residual using the preceding
parameter estimate and the midpoint field observed after the interval.
This checks compatibility on new observations before updating; it does not
forecast the next flow state or form an independent random test split.
Simulator partials may be used by the evaluator; they are
not identification targets. With noisy measurements or finite-difference error,
the relevant observation is a reproducible residual floor, not a claim that
the energy must be exactly zero.

Suppose the calibrated implicit relation has the regular readout
\begin{equation}
 \widehat A=\widehat g(X,r),\qquad y=\ell(X).
 \label{eq:afc_forward_readout}
\end{equation}
For a prescribed reference $y_d(t)$, let $e=y-y_d$ and set
$\dot y_{\rm cmd}=\dot y_d-ke$, with $k>0$. The controller completes the
output-rate relation using the measured flow state:
\begin{equation}
 r^\star\in\arg\min_{r\in\mathcal U}
 \left\{\left|D\ell(X)\widehat g(X,r)-\dot y_{\rm cmd}\right|^2
 +\lambda_r r^2\right\}.
 \label{eq:afc_control}
\end{equation}
Here $\mathcal U=[-2,2]$ is the implemented amplitude interval; an actuator
rate limit is not imposed. Applying $D\ell(X)$ to the implicit flow
relation gives the output-rate constraint
\begin{equation}
 \dot y-\widehat d(X)-\widehat g_u(X)r=0,
 \label{eq:afc_output_constraint}
\end{equation}
where $\widehat d(X)=D\ell(X)[-\widehat c_aN_h(X)
+\widehat\nu L_h(X)+K_h(X)]$ and
$\widehat g_u(X)=D\ell(X)\widehat\beta B_h$.
When $\lambda_r=0$, an unconstrained minimum-norm completion is
\begin{equation}
 r_{\rm unc}=\widehat g_u(X)^\dagger
 \left(\dot y_d-ke-\widehat d(X)\right).
 \label{eq:afc_pinv}
\end{equation}
For one input, $\widehat g_u(X)\ne0$ is required. Under an exact model,
$\lambda_r=0$, an attainable reference, and inactive constraints,
the ideal continuous-time relation is
$\dot e=-ke$. Sampling, approximation error, and active constraints leave a
completion residual. Output tracking alone does not establish stability of
all internal flow modes.

\section{Executable protocol}
\paragraph{Geometry, resolution, and sensing.}
The domain is a periodic $6\times3$ rectangle with a fringe relaxation
layer; this approximates a finite uniform-flow region. A unit-chord,
closed-trailing-edge NACA~0012 profile is centered at $(2,1.5)$ and rotated
by $12^\circ$. The default grid has $128\times64$ points, with a strict
two-thirds Fourier cutoff, penalty time $\kappa^{-1}=0.04$, and a mask
smoothing width of $0.65$ grid cells. The nominal Reynolds number is 100.
The independent flow integrator uses three-stage SSP Runge--Kutta with
$\Delta t=0.004$; fields are observed every $h_s=0.02$.
These are noiseless full-field simulation observations. Pressure, lift,
and drag are not identification channels or reported aerodynamic outputs.
Before spin-up, the field is $(1,0)$ plus a seeded solenoidal perturbation
with grid RMS $0.002$, projected into the retained Fourier space.
The complete fringe profile and other numerical choices are recorded in
\path{configs/airfoil.json}.

\paragraph{Actuator and desired output.}
The primary vertical-force Gaussian is centered at $(2.9,1.6)$, with
standard deviations $(0.30,0.22)$. An opposite lobe centered at $(2.9,0.65)$
has the same widths and is scaled so the discrete total force is zero;
both profiles are multiplied by $1-\chi$. This is a volume-force actuator,
not surface blowing or suction. The output $y=CX$ is the vertical velocity
averaged with a normalized Gaussian centered at $(3.0,1.6)$, of widths
$(0.32,0.24)$, again masked by $1-\chi$. Its input coupling is the computed
scalar $CB_h$, required to be nonzero. For control time $t\geq0$,
\begin{equation}
 y_d(t)=y_0e^{-t/1.5}+0.04\sin(2\pi t/4),
 \qquad y_0=CX(0).
\end{equation}
With $\widehat d=C[-\widehat c_aN_h(X)+\widehat\nu L_h(X)+K_h(X)]$,
$\widehat g=\widehat\beta CB_h$, $k=3$, and $\lambda_r=10^{-4}$,
the actual applied readout is
\begin{equation}
 r=\operatorname{clip}_{[-2,2]}
 \left(\frac{\widehat g\,[\dot y_d-3(y-y_d)-\widehat d]}
 {\widehat g^2+10^{-4}}\right).
 \label{eq:afc_applied_control}
\end{equation}
It is computed from the present measured field and the preceding parameter
estimate, then held during the next observation interval.

\paragraph{Causal differential learning.}
For each completed interval, form the midpoint field
$\overline X_n=(X_{n+1}+X_n)/2$ and sampled rate
$A_n=(X_{n+1}-X_n)/h_s$. With $\overline z_n=(\overline X_n,A_n,r_n)$,
the numerical incompatibility is
\begin{equation}
 \xi_n^\Delta(\theta)=
 \frac{f_\theta(\overline z_n)-f_\theta(\overline z_{n-1})}{h_s}.
 \label{eq:afc_secant}
\end{equation}
The same candidate $\theta$ is used at both endpoints. Complete nonlinear
features are differenced; their products are not differentiated by holding
one factor fixed. The update is causal because both intervals have finished
before their information is assimilated. It approximates differential
compatibility and includes sampling error.

Starting from $E_0=0$, the sampled memory and update are
\begin{align}
 E_n(\theta)&=\rho E_{n-1}(\theta)
  +\frac{1-\rho}{2}\|\xi_n^\Delta(\theta)\|_h^2,
  \qquad \rho=e^{-h_s/0.8},\\
 \theta_n^\star&\in\arg\min_{\theta\in\Theta}
 \left\{E_n(\theta)+\frac{\lambda_n}{2}
       \|D^{-1}(\theta-\theta_0)\|^2\right\},\\
 \widehat\theta_n&=0.75\widehat\theta_{n-1}+0.25\theta_n^\star,
 \qquad D=\operatorname{diag}(1,0.02,1).
\end{align}
Here $\theta_0=(0.7,0.018,0.7)$ and $\lambda_n$ is $10^{-7}$ times
the larger of $10^{-10}$ and one-third the trace of the scaled
correlation matrix. The implementation enumerates the 27 free/lower/upper
active sets of this three-parameter box problem. Relaxing toward the
feasible minimizer decreases the frozen regularized objective; this does
not imply that the online unregularized energy decreases at every sample.

\paragraph{Episodes and comparisons.}
After three time units of unactuated spin-up, three units of identification
use the bounded input
\begin{equation}
 r(t)=0.9\sum_{j=1}^3a_j\sin(2\pi f_jt+\varphi_j),
 \quad a=(0.5,0.3,0.2),\quad f=(0.45,1.15,2.35).
\end{equation}
The multisine is evaluated at $t_n=nh_s$ and held during the next observation
interval. Its phases are drawn uniformly using seed $11$ for the recorded
seed-10 run; the initial flow perturbation uses seed~10. All control branches
start from the same resulting field and estimate. The accumulated
identification memory is retained, but the preceding secant pair is cleared.
The first control interval therefore supplies an initial endpoint and
learning resumes after the second completed interval. The branches comprise no control,
a controller with the post-identification estimate frozen, and the same
controller with continuing learning. Each branch runs for six time units.
The plant actuator efficacy changes from $\beta=1$ to $0.55$ at control
time $t=2.8$; neither controller receives this change as an identification
label. For varying coefficients, differentiating the physical relation
adds a parameter-change term. On smooth portions of the physical trajectory,
$F_{\theta_{\rm true}(t)}(z)\dot z=\dot\beta B_hr$.
At the scheduled gain jump, the executable uses finite secants; a secant
spanning the change contains a parameter-change contribution and is not
expected to have zero compatibility. The fixed-coefficient relation is
recovered separately within each regime.

\paragraph{Diagnostics and validation scope.}
The saved four-panel figure reports desired and achieved output, actual
commands, pre-update secant incompatibility, and partial errors. The latter
include a relative error of the learned operator applied to one fixed smooth solenoidal perturbation
with $\delta A=0$, $\delta r=1$. This is a directional diagnostic, not an
operator-norm bound. Its denominator includes known damping. Accordingly,
the traces also report $|\widehat\beta-\beta|/|\beta|$, and a separate plot
shows all three coefficients. Final flow fields and vorticity plots show
the state being controlled. The next-interval level residual checks the
imposed calibration separately from differential compatibility.

\paragraph{Recorded prototype run.}
Table~\ref{tab:afc_recorded} reports the actual default run with seed~10;
no multi-seed uncertainty estimate is implied. Output errors use the
300 recorded post-step control samples over six time units. The continued
learner ended at
$\widehat\theta=(0.998074,0.00998102,0.555842)$, against evaluation-only
final truth $(1,0.01,0.55)$. Its final actuator-partial relative error was
$0.01062$. The median pre-update incompatibility over the final fifth of
the control episode was $5.12\times10^{-7}$ for the continued learner and
$1.62\times10^{-4}$ for the frozen estimate. The no-control branch had
very small incompatibility despite its inaccurate actuator partial, since
the input and its variations disappeared. This directly illustrates why
energy alone is not a sufficient identification diagnostic.

\begin{table}[ht]
\centering
\caption{New seed-10 numerical prototype results, in nondimensional units.
These are not original manuscript results or validated aerodynamic data.}
\label{tab:afc_recorded}
\begin{tabular}{lrrr}
\toprule
Method & Output RMS error & Input RMS & Saturation fraction\\
\midrule
No control & $0.038734$ & $0$ & $0$\\
Frozen after identification & $0.004697$ & $0.675056$ & $0$\\
Continued learning & $0.001216$ & $0.729696$ & $0$\\
\bottomrule
\end{tabular}
\end{table}

The supplied records are a numerical implementation demonstration. The
thin immersed profile has only a few grid points across its thickness;
grid, time-step, mask, penalty, and domain convergence have not been
established. The final continued-learning field has divergence RMS about
$9.3\times10^{-15}$, but its mask-weighted speed
$[\sum\chi|X|^2/\sum\chi]^{1/2}$ is approximately $0.189$; small discrete
divergence therefore does not establish accurate wall enforcement.
A surface-jet lift/drag experiment requires a resolved wall
actuator, appropriate far-field and wall conditions, stress-based force
evaluation, and independent flow validation. The present results do not
establish aerodynamic improvement, sparse-sensor control, noise robustness,
or global closed-loop stability.

\bibliographystyle{unsrt}
\bibliography{active_flow_control}
\end{document}
